\chapter{2005年上海卷（秋理）}
\section{填空题}
\begin{problem}
函数 \( f(x)=\log_{4}(x+1) \) 的反函数 \( f^{-1}(x)= \)\fillinblank{}.
\end{problem}

\begin{problem}
方程 \( 4^{x} + 2^{x} - 2 = 0 \) 的解是\fillinblank{}.
\end{problem}

\begin{problem}
直角坐标平面 \(xOy\) 中，若定点 \(A(1,2)\) 与动点 \(P(x,y)\) 满足 \(\overrightarrow{OP} \cdot \overrightarrow{OA} = 4\)，则点 \(P\) 的轨迹方程是\fillinblank{}.
\end{problem}

\begin{problem}
在 \((x - a)^10\) 的展开式中，\(x^7\) 的系数是 15，则实数 \(a = \)\fillinblank{}.
\end{problem}

\begin{problem}
若双曲线的渐近线方程为 \( y = \pm 3x \)，它的一个焦点是 \( (\sqrt{10}, 0) \)，则双曲线的方程是\fillinblank{}.
\end{problem}

\begin{problem}
将参数方程 \( \begin{cases}x=1+2\cos\theta,\\ y=2\sin\theta,\end{cases} \) （\( \theta \) 为参数）化为普通方程，所得方程是\fillinblank{}.
\end{problem}

\begin{problem}
计算： \( \lim\limits_{x\to\infty}\frac{3^{x+1}-2^{x}}{3^{x}+2^{x-1}}= \)\fillinblank{}.
\end{problem}

\begin{problem}
某班有 50 名学生. 其中 15 人选修 \(A\) 课程，另外 35 人选修 \(B\) 课程. 从班级中任选两名学生，他们是选修不同课程的学生的概率是\fillinblank{}. （结果用分数表示）
\end{problem}

\begin{problem}
在 \(\triangle ABC\) 中，若 \(A = 120^{\circ}\)，\(AB = 5\)，\(BC = 7\)，则 \(\triangle ABC\) 的面积 \(S =\)\fillinblank{}.
\end{problem}

\begin{problem}
函数 \(f(x) = \sin x + 2|\sin x|\)，\(x \in [0, 2\pi]\) 的图象与直线 \( y = k \) 有且仅有两个不同的交点，则 \( k \) 的取值范围是\fillinblank{}.
\end{problem}

\begin{problem}
有两个相同的直三棱柱，高为 \(\frac{2}{a}\)，底面三角形的三边长分别为 \(3a\)，\(4a\)，\(5a (a > 0)\). 用它们拼成一个三棱柱或四棱柱，在所有可能的情形中，全面积最小的是一个四棱柱，则 \(a\) 的取值范围是\fillinblank{}.

\bitmapfigure[max width=0.62\linewidth,max totalheight=3.2cm]{img/2005/shanghai_science/q11_fig1.png}
\end{problem}

\begin{problem}
用 \(n\) 个不同的实数 \(a_1\)，\(a_2\)，\(\cdots\)，\(a_n\) 可得到 \(n!\) 个不同的排列，每个排列为一行写成一个 \(n!\) 行的数阵. 对第 \(i\) 行 \(a_{i1}\)，\(a_{i2}\)，\(\cdots\)，\(a_{in}\)，记 \(b_i = -a_{i1} + 2a_{i2} - 3a_{i3} + \cdots + (-1)^n na_{in}\)，\(i = 1\)，\(2\)，\(3\)，\(\cdots\)，\(n!\). 例如：用 \(1\)，\(2\)，\(3\) 可得数阵如图，由于此数阵中每一列各数之和都是 12，所以，\(b_1 + b_2 + \cdots + b_6 = -12 + 2 \times 12 - 3 \times 12 = -24\). 那么，在用 \(1\)，\(2\)，\(3\)，\(4\)，\(5\) 形成的数阵中，\(b_1 + b_2 + \cdots + b_{120} =\fillinblank{}.\)

\begin{center}
\begin{tabular}{ccc}
1 & 2 & 3 \\
1 & 3 & 2 \\
2 & 1 & 3 \\
2 & 3 & 1 \\
3 & 1 & 2 \\
3 & 2 & 1
\end{tabular}
\end{center}
\end{problem}

\section{单选题}
\begin{problem}
若函数 \( f(x) = \frac{1}{2^x + 1} \)，则该函数在 \((- \infty, +\infty)\) 上是（\quad）

\choices
    {单调递减无最小值}
    {单调递减有最小值}
    {单调递增无最大值}
    {单调递增有最大值}
\end{problem}

\begin{problem}
已知集合 \(M = \{x \mid |x - 1| \leqslant 2, x \in \R\}\)，\(P = \left\{x \mid \frac{5}{x + 1} \geqslant 1, x \in \Z\right\}\)，则 \(M \cap P\) 等于（\quad）

\choices
    {\(\{x\mid 0 < x\leqslant 3,x\in \Z\}\)}
    {\(\{x\mid 0\leqslant x\leqslant 3,x\in \Z\}\)}
    {\(\{x\mid -1\leqslant x\leqslant 0,x\in \Z\}\)}
    {\(\{x\mid -1\leqslant x < 0,x\in \Z\}\)}
\end{problem}

\begin{problem}
过抛物线 \( y^{2}=4x \) 的焦点作一条直线与抛物线相交于 \(A\)，\(B\) 两点，它们的横坐标之和等于 5，则这样的直线（\quad）

\choices
    {有且仅有一条}
    {有且仅有两条}
    {有无穷多条}
    {不存在}
\end{problem}

\begin{problem}
设定义域为 \(\R\) 的函数 \(f(x) = \begin{cases} |\lg |x - 1||, & x \neq 1, \\ 0, & x = 1, \end{cases}\) 则关于 \(x\) 的方程 \(f^2(x) + bf(x) + c = 0\) 有 7 个不同实数解的充要条件是（\quad）

\choices
    {\(b < 0\) 且 \(c > 0\)}
    {\(b > 0\) 且 \(c < 0\)}
    {\(b < 0\) 且 \(c = 0\)}
    {\(b \geqslant 0\) 且 \(c = 0\)}
\end{problem}

\section{解答题}
\begin{problem}
如图，已知直四棱柱 \(ABCD - A_{1}B_{1}C_{1}D_{1}\) 中，\(AA_{1} = 2\). 底面 \(ABCD\) 是直角梯形，\(\angle A\) 是直角，\(AB \parallel CD\)，\(AB = 4\)，\(AD = 2\)，\(DC = 1\)，求异面直线 \(BC_{1}\) 与 \(DC\) 所成角的大小. （结果用反三角函数值表示）

\bitmapfigure[max width=0.42\linewidth,max totalheight=5.0cm]{img/2005/shanghai_science/q17_fig1.png}
\end{problem}

\begin{problem}
证明：在复数范围内，方程 \(|z|^2 + (1 - \i)\overline{z} - (1 + \i)z = \frac{5 - 5\i}{2 + \i}\) （\(i\) 为虚数单位）无解.
\end{problem}

\begin{problem}
如图，点 \(A\)，\(B\) 分别是椭圆 \( \frac{x^{2}}{36} + \frac{y^{2}}{20} = 1 \) 长轴的左，右端点，点 \(F\) 是椭圆的右焦点，点 \(P\) 在椭圆上，且位于 \(x\) 轴上方，\(PA \perp PF\).

\begin{enumerate}
\item 求点 \(P\) 的坐标；
\item 设 \(M\) 是椭圆长轴 \(AB\) 上的一点，\(M\) 到直线 \(AP\) 的距离等于 \(|MB|\)，求椭圆上的点到点 \(M\) 的距离 \(d\) 的最小值.
\end{enumerate}

\bitmapfigure[max width=0.42\linewidth,max totalheight=5.0cm]{img/2005/shanghai_science/q19_fig1.png}
\end{problem}

\begin{problem}
假设某市 2004 年新建住房面积 400 万平方米，其中有 250 万平方米是中低价房. 预计在今后的若干年内，该市每年新建住房面积平均比上一年增长 8\%. 另外，每年新建住房中，中低价房的面积均比上一年增加 50 万平方米. 那么，到哪一年底，

\begin{enumerate}
\item 该市历年所建中低价房的累计面积（以 2004 年为累计的第一年）将首次不少于 4750 万平方米?
\item 当年建造的中低价房的面积占该年建造住房面积的比例首次大于 85\%?
\end{enumerate}
\end{problem}

\begin{problem}
对定义域是 \(D_f\)，\(D_g\) 的函数 \(y = f(x)\)，\(y = g(x)\)，规定：函数
\(h(x) = \begin{cases}
f(x)g(x), & \text{当 } x \in D_f \text{且} x \in D_g, \\
f(x), & \text{当 } x \in D_f \text{且} x \notin D_g, \\
g(x), & \text{当 } x \notin D_f \text{且} x \in D_g.
\end{cases}\)

\begin{enumerate}
\item 若函数 \(f(x) = \frac{1}{x - 1}\)，\(g(x) = x^2\)，写出函数 \(h(x)\) 的解析式；
\item 求问题（1）中函数 \(h(x)\) 的值域；
\item 若 \(g(x) = f(x + \alpha)\)，其中 \(\alpha\) 是常数，且 \(\alpha \in [0, \pi]\)，请设计一个定义域为 \(\R\) 的函数 \(y = f(x)\)，及一个 \(\alpha\) 的值，使得 \(h(x) = \cos 4x\)，并予以证明.
\end{enumerate}
\end{problem}

\begin{problem}
在直角坐标平面中，已知点 \(P_{1}(1,2)\)，\(P_{2}(2,2^{2})\)，\(P_{3}(3,2^{3})\)，\(\cdots\)，\(P_{n}(n,2^{n})\)，其中 \(n\) 是正整数，对平面上任一点 \(A_{0}\)，记 \(A_{1}\) 为 \(A_{0}\) 关于点 \(P_{1}\) 的对称点，\(A_{2}\) 为 \(A_{1}\) 关于点 \(P_{2}\) 的对称点，\(\cdots\)，\(A_{n}\) 为 \(A_{n-1}\) 关于点 \(P_{n}\) 的对称点.

\begin{enumerate}
\item 求向量 \(\overrightarrow{A_{0}A_{2}}\) 的坐标；
\item 当点 \(A_{0}\) 在曲线 \(C\) 上移动时，点 \(A_{2}\) 的轨迹是函数 \(y = f(x)\) 的图象，其中 \(f(x)\) 是以 3 为周期的周期函数，且当 \(x \in (0,3]\) 时，\(f(x) = \lg x\). 求以曲线 \(C\) 为图象的函数在 \((1,4]\) 上的解析式；
\item 对任意偶数 \(n\)，用 \(n\) 表示向量 \(\overrightarrow{A_{0}A_{n}}\) 的坐标.
\end{enumerate}
\end{problem}
